%! AP Statistics — Unit 2, Lesson 2.6 — Guided Notes — Answer Key
\documentclass[10pt]{article}
\usepackage{apstats-article}
\usepackage{apstats-key}

\begin{document}

\pageheader{Unit 2, Lesson 2.6}{Guided Notes --- Answer Key}
\namedateperiod

\vspace{0.12in}

% ── OBJECTIVE ────────────────────────────────────────────────────────────────
\begin{objectivebox}
\small By the end of this lesson, I will be able to\dots\\[4pt]
\ansline{write the LSRL equation using variable names and interpret the slope and $y$-intercept in context.}\\[1pt]
\ansline{use the regression equation to make predictions and distinguish interpolation from extrapolation.}\\[1pt]
\ansline{explain when the $y$-intercept is or is not meaningful in context.}
\end{objectivebox}

\vspace{0.1in}

% ── VOCABULARY ───────────────────────────────────────────────────────────────
\begin{vocabbox}
\termblanklong{Regression line (LSRL)}
\ansline{The line $\hat{y} = a + bx$ that best summarizes a linear association by minimizing the sum of squared residuals.}
\termblanklong{Slope ($b$)}
\ansline{The predicted change in $\hat{y}$ for each one-unit increase in $x$; includes sign (direction) and units.}
\termblanklong{$y$-intercept ($a$)}
\ansline{The predicted value of $\hat{y}$ when $x = 0$; not always meaningful in context.}
\termblanklong{$\hat{y}$ (``$y$-hat'' / predicted value)}
\ansline{The value of $y$ estimated from the regression equation for a given $x$; individual values will differ from $\hat{y}$.}
\termblanklong{Interpolation}
\ansline{Predicting $\hat{y}$ for an $x$ value within the observed data range; generally reliable.}
\termblanklong{Extrapolation}
\ansline{Predicting $\hat{y}$ for an $x$ value outside the observed data range; unreliable; should be flagged.}
\end{vocabbox}

\vspace{0.1in}

% ── HOOK ─────────────────────────────────────────────────────────────────────
\begin{hookbox}
\small\textbf{Opening Question:} Your teacher shows a scatterplot of height (inches) vs.\
shoe size for 30 students and asks: ``If a student is 68 inches tall, what shoe size would
you predict?''

\vspace{6pt}
Write your estimate: Predicted shoe size $\approx$ \ans{9--10 (accept reasonable estimates)} \quad How did you decide?\\[8pt]
\ansline{Answers vary; students should describe reading from or sketching a line through the cloud of points.}

\vspace{4pt}
\textbf{Key idea:} Today we formalize this process with an equation --- the \textbf{regression line}.
\end{hookbox}

\vspace{0.1in}

% ── STEP 1 ───────────────────────────────────────────────────────────────────
\begin{notesbox}{Step 1 --- The Regression Equation (DAT-1.E)}
\small

\noindent The equation of the \textbf{least-squares regression line (LSRL)} is:
\[
  \hat{y} = a + bx
\]
where \quad $a = $ \ans{the $y$-intercept} \quad and \quad $b = $ \ans{the slope}.

\vspace{6pt}
\textbf{AP rule:} Always write the equation using the \ans{names (variable names)} of the variables,
not generic $x$ and $y$.

\vspace{8pt}
\textbf{Class example} --- age of a used car ($x$, years) vs.\ resale price ($\hat{y}$, \$1000s):

\vspace{4pt}
Technology output gives $a = 28.4$ and $b = -1.85$. Write the equation using variable names:

\vspace{6pt}
\hspace{1.5cm} \ans{$\widehat{\text{price}} = 28.4 - 1.85(\text{age})$}

\end{notesbox}

\vspace{0.1in}

% ── STEP 2 ───────────────────────────────────────────────────────────────────
\begin{notesbox}{Step 2 --- Interpreting the Slope (DAT-1.E)}
\small

\noindent The \textbf{slope} $b$ is the predicted \ans{change} in $\hat{y}$ for each
\ans{one-unit} increase in $x$.

\vspace{6pt}
\textbf{AP template:} ``For each additional \ans{[one unit of $x$ in context]}, the predicted
\ans{[y-variable]} [increases / decreases] by approximately \ans{[|b| with units]}.''

\vspace{8pt}
\textbf{Interpret the slope} ($b = -1.85$) for the car example:
\vspace{4pt}\\
\ansline{For each additional year of age, the predicted resale price of a used car decreases by approximately \$1{,}850.}

\vspace{6pt}
\textbf{Every AP slope interpretation needs all three:}
\begin{enumerate}[leftmargin=1.8em, itemsep=2pt]
  \item \ans{``For each additional year of age''} \hfill (one-unit change in $x$, in context)
  \item \ans{``the predicted price \emph{decreases}''} \hfill (direction: increases or decreases)
  \item \ans{``by approximately \$1{,}850''} \hfill (magnitude with units)
\end{enumerate}

\end{notesbox}

\vspace{0.1in}

% ── STEP 3 ───────────────────────────────────────────────────────────────────
\begin{notesbox}{Step 3 --- Interpreting the $y$-Intercept (DAT-1.E)}
\small

\noindent The \textbf{$y$-intercept} $a$ is the predicted value of $\hat{y}$ when $x = $ \ans{0}.

\vspace{6pt}
\textbf{AP template:} ``When [x-variable] $= 0$, the predicted [y-variable] is approximately
\ans{[intercept value with units]}.''

\vspace{8pt}
\textbf{Interpret the $y$-intercept} ($a = 28.4$) for the car example:
\vspace{4pt}\\
\ansline{When a car's age is 0 years (brand new), the predicted resale price is approximately \$28{,}400.}

\vspace{6pt}
\textbf{Is this meaningful in context?} \ans{Yes} \quad \textbf{Why or why not?}\\[6pt]
\ansline{A brand-new car (age = 0) is plausible, so the intercept has a reasonable interpretation here.}

\vspace{4pt}
\begin{tcolorbox}[colback=redbg, colframe=redacc, boxrule=0.8pt, arc=2mm,
    left=3mm, right=3mm, top=2mm, bottom=2mm]
\small\textbf{Watch out:} If $x = 0$ is \ans{outside} the data range or physically
\ans{impossible}, state that the $y$-intercept is \textbf{not meaningful in context}.
\end{tcolorbox}

\end{notesbox}

\vspace{0.1in}

% ── STEP 4 ───────────────────────────────────────────────────────────────────
\begin{notesbox}{Step 4 --- Predictions: Interpolation vs.\ Extrapolation (DAT-1.F)}
\small

\noindent To \textbf{predict} $\hat{y}$, substitute an $x$ value into the regression equation.

\vspace{6pt}
\textbf{Predict the price of a 5-year-old car:}

\vspace{4pt}
\hspace{1.5cm} $\widehat{\text{price}} = 28.4 - 1.85(\ans{5})
  = \ans{19.15}$ thousand dollars

\vspace{8pt}
\renewcommand{\arraystretch}{1.7}
\begin{tabularx}{\linewidth}{@{} >{\bfseries}p{3.2cm} X @{}}
\rowcolor{sky} \textbf{Term} & \textbf{Definition} \\
Interpolation &
  $x$ is \ans{within} the data range; prediction is generally \ans{reliable}. \\
Extrapolation &
  $x$ is \ans{outside} the data range; prediction is \ans{unreliable} and should be flagged. \\
\end{tabularx}

\vspace{8pt}
\textbf{The car data ranged from age 1 to 10 years.}

\vspace{4pt}
Predict the price of a \textbf{20-year-old} car:
\hspace{0.5cm} $\widehat{\text{price}} = 28.4 - 1.85(20) = \ans{-8.6}$ thousand dollars

\vspace{4pt}
Is this interpolation or extrapolation? \ans{Extrapolation} \quad Why is this prediction unreliable?\\[6pt]
\ansline{Age 20 is far outside the data range (1--10 years), so the linear pattern may not continue. The
negative predicted price is nonsensical, illustrating the danger of extrapolation.}

\end{notesbox}

\vspace{0.1in}

% ── STEP 5 ───────────────────────────────────────────────────────────────────
\begin{notesbox}{Step 5 --- The LSRL Passes Through $(\bar{x},\, \bar{y})$}
\small

\noindent The regression line always passes through the point \ans{$(\bar{x},\, \bar{y})$} --- the mean of $x$
and the mean of $y$.

\vspace{6pt}
\textbf{Verify for the car data:} $\bar{x} = 5.5$ years, \; $\bar{y} = \ans{18.225}$ thousand dollars.

\vspace{4pt}
\hspace{1.5cm} $\widehat{\text{price}} = 28.4 - 1.85(5.5) = \ans{18.225}$ \quad Matches $\bar{y}$? \ans{Yes}

\end{notesbox}

\vspace{0.1in}

% ── CONNECTIONS ───────────────────────────────────────────────────────────────
\begin{spiralbox}
\small
\begin{multicols}{2}

\textbf{How this connects:}
\begin{itemize}[itemsep=3pt]
  \item Lesson 2.5: $r$ describes direction and strength; slope $b$ gives the rate of change.
  \item Lesson 2.7: each point's \textit{residual} $= y - \hat{y}$ measures prediction error.
  \item Lesson 2.8: the LSRL minimizes the sum of squared residuals.
\end{itemize}

\columnbreak

\textbf{Reflection:}\\
Why is extrapolation unreliable? Give your own example of a prediction that would involve
extrapolation.\\[2pt]
\ansline{The linear pattern was built from data in a specific range; outside that range, the relationship may
curve, reverse, or break down. Example: predicting a person's height at age 60 using a regression of
height on age built from data on children ages 2--12.}

\end{multicols}
\end{spiralbox}

\vspace{0.1in}

\begin{notesbox}{Additional Notes}
\vspace{0pt}
\writeline\\\writeline\\\writeline\\\writeline\\\writeline\\\writeline
\end{notesbox}

\begin{teachernote}
Step 4 verification: $28.4 - 1.85(20) = 28.4 - 37.0 = -8.6$. This absurd result (negative price)
is a powerful hook for why extrapolation is unreliable --- use it. Step 5: $28.4 - 1.85(5.5)
= 28.4 - 10.175 = 18.225$, confirming the line passes through $(\bar{x}, \bar{y})$.
\end{teachernote}

\end{document}
